% !TeX root = ../se200_identity_regimes.tex

\subsection{Refinement of Applicability Contexts}
\label{subsec:ctx_split}

This subsection shows that applicability contexts do not define a single
regime profile under the transformation basis, and must be refined.

%--------------------------------------------
\subsubsection{Initial CTX Profile}

\begin{definition}[Provisional CTX profile]
A provisional CTX profile is a regime profile whose identity carrier is
an applicability scope or condition, i.e., a structure that determines
where, to whom, or under what conditions authority applies.
\end{definition}


Under this provisional formulation, identity is intended to track persistence
of applicability independently of internal structure.


%--------------------------------------------
\subsubsection{Transformation Behavior}

We examine the classification of transformations under the provisional CTX profile.

\begin{itemize}
\item $\mathrm{RE}$: IGN
\item $\mathrm{AN}$: IGN
\item $\mathrm{RF}$: PRS or BRK
\item $\mathrm{AD}$: PRS or BRK
\item $\mathrm{RC}$: PRS or BRK
\item $\mathrm{RA}$: IGN or BRK
\item $\mathrm{SU}$: BRK
\item $\mathrm{BF}$: PRS or BRK
\item $\mathrm{PV}$: IGN
\item $\mathrm{SE}$: not applicable
      ($\mathrm{Appl}_P(\mathrm{SE}) = 0$;
      normative structures do not undergo state evolution).
\end{itemize}


The classification of decomposition is indeterminate under the provisional CTX profile:
it is identity-preserving when applicability extension is unchanged,
but identity-breaking when internal structure is changed.
This indeterminacy is the source of split pressure.


%--------------------------------------------
\subsubsection{Split Pressure}

\begin{proposition}[CTX split pressure]
\label{prop:ctx_split_pressure}
The provisional CTX profile is under split pressure with respect to
the decomposition transformation $\mathrm{AD}$.
\end{proposition}

\begin{proof}[Proof of Proposition~\ref{prop:ctx_split_pressure}]
Let $x \in \mathcal{R}$ be a context representation whose identity is
determined by applicability extension:
$x$ represents a jurisdiction defined as the set of cases
falling under a given authority,
independently of how that jurisdiction is internally organized.
Let $\tau$ be a decomposition transformation producing subcontexts
whose union covers the same set of cases as $x$, and
let $y := \tau(x)$.
Since the applicability extension of
$y$ equals that of $x$,
identity-preservation is required:
$\mathrm{Class}_P(\mathrm{AD}) = \mathrm{PRS}$.

Let $x' \in \mathcal{R}$ be a context representation whose identity is
determined by internal structure: $x'$ represents a partition of cases
into named subregions, where the partition itself constitutes the context.
Let $\tau'$ be a decomposition transformation that subdivides $x'$
into finer subregions, producing
$y' := \tau'(x')$ with different internal organization.
Since the structure of $y'$ differs from that of $x'$,
identity-breaking is required: $\mathrm{Class}_P(\mathrm{AD}) = \mathrm{BRK}$.

Both $x$ and $x'$ are admissible instances of the provisional CTX profile,
as both are applicability scopes independent of normative or causal content.
The requirements $\mathrm{PRS}$ and $\mathrm{BRK}$ are contradictory.
Therefore the provisional CTX profile is under split pressure
with respect to $\mathrm{AD}$.
\end{proof}

[Informally]
After decomposition, the applicability can remain the same even when the structure does not.


%--------------------------------------------
\subsubsection{Separated Profiles}

We therefore define two distinct regime profiles.

\begin{definition}[CTX-E profile]
The extensional applicability profile $\mathrm{CTX\mbox{-}E}$ is defined by:
\begin{itemize}
\item carrier: applicability scope or condition,
\item identity basis: persistence of the same applicability extension,
\item transformation classification:
\begin{itemize}
  \item $\mathrm{RE}$: IGN
  \item $\mathrm{AN}$: IGN
  \item $\mathrm{RF}$: PRS
  \item $\mathrm{AD}$: PRS
  \item $\mathrm{RC}$: PRS \quad (re-contextualization preserves
        applicability extension by definition of CTX-E)
  \item $\mathrm{RA}$: IGN \quad (reassignment does not alter
        applicability extension)
  \item $\mathrm{SU}$: BRK
  \item $\mathrm{BF}$: PRS
  \item $\mathrm{PV}$: IGN
  \item $\mathrm{SE}$: not applicable
      ($\mathrm{Appl}_P(\mathrm{SE}) = 0$;
      normative structures do not undergo state evolution).
\end{itemize}
\end{itemize}
\end{definition}

\begin{definition}[CTX-S profile]
The structured applicability profile $\mathrm{CTX\mbox{-}S}$ is defined by:
\begin{itemize}
\item carrier: applicability structure,
\item identity basis: persistence of the same structure,
\item transformation classification:
\begin{itemize}
  \item $\mathrm{RE}$: IGN
  \item $\mathrm{AN}$: IGN
  \item $\mathrm{RF}$: BRK
  \item $\mathrm{AD}$: BRK
  \item $\mathrm{RC}$: BRK \quad (re-contextualization alters
        applicability structure)
  \item $\mathrm{RA}$: BRK \quad (reassignment alters the structural
        bearer of applicability)
  \item $\mathrm{SU}$: BRK
  \item $\mathrm{BF}$: BRK
  \item $\mathrm{PV}$: IGN
  \item $\mathrm{SE}$: not applicable
      ($\mathrm{Appl}_P(\mathrm{SE}) = 0$;
      normative structures do not undergo state evolution).
\end{itemize}
\end{itemize}
\end{definition}


The classifications of $\mathrm{RC}$ and $\mathrm{RA}$ are fully determined
by the identity basis of each profile. Under CTX-E, both are PRS or IGN
because neither alters applicability extension. Under CTX-S, both are BRK
because each alters the structural configuration that constitutes identity.

This does not induce further split pressure.
In both profiles, re-contextualization or reassignment is identity-preserving
if and only if the defining identity basis is unchanged:
applicability extension for $\mathrm{CTX\mbox{-}E}$,
and applicability structure for $\mathrm{CTX\mbox{-}S}$.

Thus the variation in classification reflects dependence on whether the
underlying invariant is preserved, rather than the presence of distinct
regime profiles.


%--------------------------------------------
\subsubsection{Non-Collapse}

\begin{proposition}[CTX-E and CTX-S do not collapse]
\label{prop:ctx_noncollapse}
The profiles $\mathrm{CTX\mbox{-}E}$ and $\mathrm{CTX\mbox{-}S}$ induce
distinct equivalence relations on $\mathcal{R}$.
\end{proposition}

\begin{proof}[Proof of Proposition~\ref{prop:ctx_noncollapse}]
Let $x \in \mathcal{R}$ be a context representation.
Let $\tau$ be a decomposition transformation producing subcontexts
whose union has the same applicability extension as $x$,
and let $y := \tau(x)$.

Under $\mathrm{CTX\mbox{-}E}$:
\[
\mathrm{Class}_{\mathrm{CTX\mbox{-}E}}(\mathrm{AD}) = \mathrm{PRS},
\]
so $\tau$ is identity-preserving and $x \sim_{\mathrm{CTX\mbox{-}E}} y$.

Under $\mathrm{CTX\mbox{-}S}$:
\[
\mathrm{Class}_{\mathrm{CTX\mbox{-}S}}(\mathrm{AD}) = \mathrm{BRK},
\]
so $\tau$ is identity-breaking and $x \not\sim_{\mathrm{CTX\mbox{-}S}} y$.

Therefore $\sim_{\mathrm{CTX\mbox{-}E}} \neq \sim_{\mathrm{CTX\mbox{-}S}}$.
\end{proof}

[Informally]
After decomposition, two representations can have the same applicability
even when their structure differs.


%--------------------------------------------
\subsubsection{Result}

\begin{theorem}[CTX refinement]
Applicability contexts must be refined into at least two distinct regime profiles:
\[
\mathrm{CTX\mbox{-}E}
\quad\text{and}\quad
\mathrm{CTX\mbox{-}S}.
\]
\end{theorem}

[Informally]
Two representations can apply to the same cases even when their internal structure differs,
but cannot be the same structure once decomposition changes that organization.

